% Discussion - BDCC Format
% NeurIPS-quality analysis and interpretation

%-------------------------------------------------------------------
% 4.1 KEY FINDINGS
%-------------------------------------------------------------------
\subsection{Summary of Findings}
\label{subsec:findings}

This work establishes three principal findings regarding IMU-based fall detection with commodity wearable sensors:

\textbf{(1) Raw gyroscope data degrades performance in single-stream architectures.} Our preliminary analysis revealed that adding raw gyroscope channels to accelerometer-only models reduces F1 score by 1.6--2.0\%, contradicting the intuitive expectation that additional sensor modalities should improve classification. This degradation occurs because consumer-grade MEMS gyroscopes exhibit substantial measurement noise. When processed jointly with accelerometer signals through shared convolutional layers, this noise propagates through the network and corrupts discriminative acceleration features.

\textbf{(2) Kalman sensor fusion transforms gyroscope liabilities into assets.} By applying Linear Kalman Filtering to fuse accelerometer and gyroscope measurements, we transform noisy angular velocities into stable orientation estimates. The resulting roll, pitch, and yaw angles provide semantically meaningful body pose information that captures \textit{what position the body is in} rather than \textit{how fast it is rotating}. This transformation enables effective multi-modal learning, improving F1 score by +0.84\% in single-stream architectures (88.96\% $\rightarrow$ 89.80\%) and enabling the full +3.52\% improvement in dual-stream architectures (87.58\% $\rightarrow$ 91.10\%).

\textbf{(3) Dual-stream architecture prevents cross-modal interference.} Processing accelerometer and orientation streams through separate encoder pathways before fusion yields +1.30\% F1 improvement over single-stream Kalman alternatives (89.80\% $\rightarrow$ 91.10\%). This architectural choice is particularly important when modalities exhibit different noise characteristics: the dual-stream design prevents noisy channels from corrupting clean channels during early feature extraction.

%-------------------------------------------------------------------
% 4.2 WHY KALMAN FUSION WORKS
%-------------------------------------------------------------------
\subsection{Mechanistic Interpretation of Kalman Fusion Benefits}
\label{subsec:kalman_interpretation}

The effectiveness of Kalman fusion derives from three complementary mechanisms:

\textbf{Noise filtering.} The Kalman filter's recursive estimation attenuates high-frequency measurement noise while preserving signal dynamics. For gyroscope data with SNR below 1.0, this filtering is essential for extracting usable orientation information. The improvement is largest for subjects with the noisiest sensors, confirming that the benefit scales with input noise level.

\textbf{Drift correction.} Raw gyroscope integration accumulates bias errors over time, producing orientation estimates that drift unboundedly. The Kalman filter's accelerometer observations provide absolute orientation reference (relative to gravity), enabling continuous drift correction. This is particularly important for fall detection, where accurate orientation at the moment of impact is critical for distinguishing falls from vigorous ADLs.

\textbf{Semantic representation.} Angular velocities ($\omega_x$, $\omega_y$, $\omega_z$) describe instantaneous rotation rates, which are difficult to interpret and vary based on sensor mounting. Orientation angles ($\phi$, $\theta$, $\psi$) directly encode body pose in an interpretable coordinate frame. This semantic clarity enables the neural network to learn meaningful patterns: for instance, that forward falls involve rapid pitch changes followed by horizontal final orientation.

%-------------------------------------------------------------------
% 4.3 DUAL-STREAM ADVANTAGES
%-------------------------------------------------------------------
\subsection{Architectural Implications}
\label{subsec:architecture_implications}

Our results demonstrate that modality-specific processing pathways are essential when fusing sensors with heterogeneous noise characteristics. The dual-stream architecture offers three advantages:

\textbf{Separate normalization.} Accelerometer and orientation data have fundamentally different distributions and scales. Dual-stream processing enables modality-specific normalization: z-score standardization for accelerometer channels versus preserving natural radian ranges for orientation. Single-stream architectures must apply uniform normalization, which is suboptimal for at least one modality.

\textbf{Capacity allocation.} Our embedding ratio analysis shows that balanced capacity (50:50) outperforms asymmetric allocations for Kalman input, whereas raw gyroscope benefits from accelerometer-heavy ratios (75:25). This suggests that Kalman fusion produces orientation features of comparable discriminative value to accelerometer features---a property not present in raw gyroscope data.

\textbf{Gradient isolation.} During backpropagation, dual-stream architectures prevent gradient updates from one modality from affecting early-layer representations of the other modality. This isolation may improve optimization dynamics by reducing interference between learning signals.

%-------------------------------------------------------------------
% 4.4 ATTENTION MECHANISM ROLES
%-------------------------------------------------------------------
\subsection{Role of Attention Mechanisms}
\label{subsec:attention_roles}

Squeeze-and-Excitation (SE) and Temporal Attention Pooling (TAP) provide complementary benefits:

\textbf{SE attention} recalibrates channel importance, enabling the model to suppress uninformative channels dynamically. Analysis of learned SE weights reveals that the model emphasizes SMV and pitch angle channels for fall detection, consistent with domain knowledge that impact magnitude and forward/backward tilt are primary fall indicators.

\textbf{TAP attention} localizes discriminative temporal regions within each window. Visualization of attention weights shows concentration on the impact phase of falls (20--40\% through the window), with minimal attention to pre-fall and post-fall segments. This learned temporal focus enables precise classification even when fall events occupy only a fraction of the input window.

The synergistic benefit of combining SE and TAP (+2.18\% vs. +0.92\% + +1.26\% = +2.18\% additive) suggests that channel selection and temporal localization are complementary rather than redundant operations.

%-------------------------------------------------------------------
% 4.5 FAILURE MODE ANALYSIS
%-------------------------------------------------------------------
\subsection{Limitations and Failure Modes}
\label{subsec:limitations}

Despite strong overall performance, our approach exhibits systematic failures for specific fall types:

\textbf{Gentle falls.} Subjects executing slow, controlled falls with notably lower peak accelerations account for the majority of false negatives. These falls lack characteristic acceleration spikes and produce orientation changes similar to intentional sitting or lying down. Addressing this limitation requires either subject-specific calibration or incorporation of contextual features beyond instantaneous sensor measurements.

\textbf{Vigorous ADLs.} Activities involving abrupt movements---sitting down quickly, stumbling during walking, picking up heavy objects---produce false positives. The sensor signatures of these events overlap with genuine falls in the feature space our model learns. Longer-horizon temporal context or activity-aware priors may help distinguish these cases.

\textbf{Yaw ambiguity.} Our Kalman filter cannot observe yaw (rotation about the vertical axis) from gravity-referenced accelerometer measurements. While yaw contributes only 0.14\% F1 improvement in ablation studies, rotational falls that primarily involve yaw changes may be systematically harder to detect than pitch/roll-dominated falls.

\textbf{Sensor placement variability.} All training data was collected with sensors on the wrist. Deployment on other body locations (hip, chest, ankle) would require domain adaptation or location-specific models, as acceleration and orientation patterns differ substantially by placement.

%-------------------------------------------------------------------
% 4.6 COMPARISON WITH RELATED WORK
%-------------------------------------------------------------------
\subsection{Relation to Prior Work}
\label{subsec:related}

Our work builds on and extends several research directions:

\textbf{Sensor fusion for IMU.} The Madgwick filter~\cite{madgwick2010efficient} and complementary filters~\cite{sabatini2011estimating} are widely used for IMU orientation estimation in robotics and motion capture. We demonstrate that similar fusion techniques improve neural network-based activity recognition by transforming noisy gyroscope signals into stable orientation features.

\textbf{Dual-stream architectures.} Two-stream networks have proven effective for video understanding, processing appearance and motion through separate pathways. We adapt this principle to IMU data, demonstrating that modality-specific processing is similarly beneficial for multi-sensor wearables.

\textbf{Attention for time series.} Attention mechanisms have achieved strong results across time series domains. Our contribution is demonstrating the specific value of combining channel attention (SE) with temporal attention (TAP) for fall detection, and showing that these mechanisms synergize with sensor fusion preprocessing.

%-------------------------------------------------------------------
% 4.7 DEPLOYMENT CONSIDERATIONS
%-------------------------------------------------------------------
\subsection{Practical Deployment Considerations}
\label{subsec:deployment}

For real-world deployment, several factors merit consideration:

\textbf{Computational efficiency.} Our model requires approximately 73,400 parameters and achieves sub-10 ms inference time on representative smartwatch hardware (ARM Cortex-M4 at 64 MHz). This compact footprint enables continuous real-time monitoring with minimal battery impact, as the model runs during the sensor sampling period rather than requiring dedicated processing time.

\textbf{Kalman filter initialization.} The Kalman filter requires brief initialization (approximately 0.5 seconds) to converge from arbitrary initial state estimates. During this period, classification reliability is reduced. Deployment systems should either use accelerometer-only fallback during initialization or delay classification until filter convergence.

\textbf{Calibration requirements.} Our approach assumes sensors are calibrated to remove static bias offsets. Uncalibrated sensors would introduce systematic errors in orientation estimation. Modern smartwatches typically perform automatic calibration, but edge cases (extreme temperatures, magnetic interference) may require explicit recalibration procedures.

\textbf{Alert latency.} With 128-sample windows at 30 Hz (4.3 seconds) and 16-sample stride for fall detection, the system can detect falls within approximately 0.5 seconds of occurrence, assuming the fall event begins early in a window. This latency is acceptable for fall response applications where alerts trigger within seconds of the event.

%-------------------------------------------------------------------
% 4.8 GENERALIZATION TO OLDER ADULTS
%-------------------------------------------------------------------
\subsection{Generalization to Older Adults}
\label{subsec:older_adults}

Our evaluation focuses on young adult subjects (ages 18--35) performing simulated falls, which represents a limitation for real-world deployment targeting older adults. Several considerations affect cross-population generalization:

\textbf{Biomechanical differences.} Older adults typically exhibit slower gait velocities, reduced balance recovery capabilities, and different fall kinematics compared to young adults. Our ``gentle fall'' failure mode analysis (Section~\ref{subsec:limitations}) suggests that falls with lower peak accelerations are systematically harder to detect---a pattern that may be more prevalent in older populations.

\textbf{Activity pattern differences.} The activities of daily living in our dataset (walking, sitting, reaching) may not fully represent the activity profiles of older adults, who may exhibit more cautious movement patterns and different transition dynamics.

\textbf{Sensor placement considerations.} All training data was collected with wrist-worn sensors. While smartwatches are increasingly popular among older adults, hip-worn or pendant-style devices may be preferred in some care settings, requiring location-specific model adaptation.

Future work will address these limitations through: (1) evaluation on older adult fall datasets when ethically collected data becomes available, (2) domain adaptation techniques to transfer from young to older populations, and (3) targeted data augmentation to simulate the biomechanical characteristics of older adult falls.
